% !TeX root = ../机器学习笔记.tex
% 由主文档通过 \input 载入；请编译主文档。
\setcounter{chapter}{9} % 保留原书章号。
\chapter{近邻方法与非参数估计}
\label{chap:ml-neighbors}

近邻方法把“相似输入具有相似输出”作为主要假设。它不先拟合一个固定维度的全局函数，而是在预测时直接利用样本附近的信息。

\section{距离、相似性与特征尺度}

Minkowski 距离为 $d_p(x,z)=(\sum_j|x_j-z_j|^p)^{1/p}$，$p\geq1$；$p=1,2$ 分别对应曼哈顿距离和欧氏距离。非零向量的余弦相似度为 $x^\top z/(\|x\|\|z\|)$，它本身不是距离。正定矩阵 $M$ 定义 Mahalanobis 距离 $\sqrt{(x-z)^\top M(x-z)}$；仅半正定时可能只得到伪度量。

距离依赖单位。把米改为毫米会使对应维度在欧氏距离中支配其他特征，因此尺度处理是建模的一部分。无关特征增多也会稀释有效邻近关系，标准化不能把无关特征变成有用特征。

\section{近邻分类与回归}

$k$ 近邻（$k$-nearest neighbors，$k$-NN）保存训练样本，查询时找出距离最近的 $k$ 个索引 $\mathcal N_k(x)$。回归预测为 $\widehat f(x)=k^{-1}\sum_{i\in\mathcal N_k(x)}y_i$，分类以邻居类别频数估计后验，再作多数投票。也可使用非负归一权重 $a_i(x)$，令预测为 $\sum_i a_i(x)y_i$。

算法需要事先规定距离并列、类别并列和零距离的处理。使用倒距离权重时，遇到重合点应直接在重合样本内预测或采用明确的平滑规则，不能直接除以零。$k$ 是验证选择的超参数，单近邻在没有冲突标签和重复点干扰时可记住训练样本，但这不代表测试最优。

\section{局部平均、平滑与一致性}

假设 $Y=m(X)+\varepsilon$，条件均值为零的噪声在邻居间独立，条件方差至多 $\sigma^2$。若 $m$ 为 $L_m$-Lipschitz，邻居距查询点至多 $r$，则给定输入位置后的预测偏差绝对值至多 $L_mr$，方差至多 $\sigma^2/k$。因此条件均方误差至多为 $L_m^2r^2+\sigma^2/k$。前一项偏好小邻域，后一项偏好多邻居。

若在固定查询点附近存在正密度，增大样本量时可同时令邻域缩小、邻居数增大。这解释 $k\to\infty$、$k/n\to0$ 的常见选择路线。上述局部计算不是一般一致性定理；一般分布、距离并列和全局可积性还需单独处理。本书只使用这里已陈述的光滑回归假设下的误差解释。

\section{核密度与核回归}

取非负且积分为一的平滑核 $\kappa:\mathbb R^d\to\mathbb R$，带宽 $h>0$，核密度估计（kernel density estimation）为
\begin{equation}
 \widehat p_h(x)=\frac1{nh^d}\sum_i\kappa((x-x_i)/h).
 \label{eq:ml-kde}
\end{equation}
变量替换可验证其积分为一。核回归则取
\[
 \widehat m_h(x)=
 \frac{\sum_i\kappa((x-x_i)/h)y_i}{\sum_i\kappa((x-x_i)/h)}
\]
且要求分母正。带宽决定局部范围，太小会有高方差，太大会抹去局部变化。此处的核是平滑权重，与第\ref{chap:ml-kernels}章的半正定核要求不同。

\section{搜索结构与维数灾难}

稠密数据上一次朴素查询需 $O(nd)$ 距离计算，另有选取邻居的成本；保存数据需 $O(nd)$ 空间。$k$-d 树等索引通过空间划分减少不必访问的点，但其收益依赖维度和数据分布，最坏情况下仍可能扫描大量样本。近似检索交换查询成本与邻居精确性。

在 $d$ 维区域内，半径按比例缩小的邻域体积约按该比例的 $d$ 次方缩小。维度升高后，获得同样局部样本数可能需要急剧增加总体样本量。这是维数灾难（curse of dimensionality）的一种表现，也说明低维有效表示为何重要。

\section{例题、练习与学习检查}

\begin{example}[邻居数量的影响]
训练输入为 $0,1,3$，响应为 $0,2,2$，查询 $x=0.2$。
\end{example}
\begin{solve}
距离为 $0.2,0.8,2.8$。单近邻预测为零，两近邻预测为一，三近邻预测为 $4/3$。扩大邻域使结果更平滑，也纳入了较远样本；仅凭一个查询无法确定哪个 $k$ 最优。
\end{solve}

习题：证明归一核密度的积分为一；推导本章局部均方误差界；把一个输入维度乘以一千，观察邻居次序是否变化；比较训练单近邻误差与留一评价误差。





